\chapter{The joint rational C2 opening: a full-parent certificate}
\label{component:p8f912583feaa}
\noindent\textbf{Component }\path{p8f912583feaa}\par
{\small\noindent Source: \path{R13/baseline/baseline/baseline/baseline/advances/c2_joint/joint_centered_packets.tex}\par}
{\small\noindent SHA256: \path{8f912583feaa7f930694a5996b34ad984e1af7b157dbfae029f41c58a4a2ef70}\par}
{\small\noindent Inventory use: JC1--JC13 complete simultaneous r,gInf source with both centered excess packets\par}
\medskip
\section[The joint rational C2 opening: a full-parent certificate]{The joint rational C2 opening: a full-parent certificate}
\label{p8f912583feaa:sec:c2-joint-centered}

\begin{theorem}[Both event openings may close]
Use the exact rational C2 source and original parent of
\path{general_single_angle.tex}.  Suppose $x>0$ is infinitesimal,
$k,b$ are positive bounded units, $0<r<1$, $0<h<b$, and
$g=b-1+r+xh>0$.  Suppose $r,g$ are infinitesimal, $n\geq0$, and
the source is physical.  Then the entire original parent is tensor
SOS at every independent real or complex projective testing state
and every independent projective pair of testing angles.  No lower
bound on either $h$ or $b-h$ is required.
\end{theorem}

Together with the three-opening and last-opening theorems this
removes the condition that at least one of $r,g$ be a unit in the
bounded-$k,b$ rational C2 source family.  It does not classify all
physical one-weak C2 sources, nor admit a nonunit $k$ or unbounded $b$.
All tensor-ring conventions and the strict bounded tensor Gram lemma
are those of the authoritative package.  In particular, a root vector
need not be linearly independent, and its entries may be unbounded.

\paragraph{Source bounds, including both complementary openings.}
Write $\rho=\sqrt{x}$, $\eta=h/b$, $\kappa=b/k$, $K=k_0>0$,
and $D_K=1+K^2$.  Since $r,g$ are infinitesimal, $b_0=1$.
The source formulas give
\[
 (nr)_0=\frac1{4D_K},\qquad (mg)_0=\frac14,
 \qquad \Delta_0=2D_K.
\]
Thus $n,m$ are positive unbounded pure coefficient quantities,
$nr,mg$ are bounded units, and $nr^2,mg^2$ are infinitesimal.
In particular $m-1>0$.  The existing opening cancellation identities,
with $T_C=xC(b-h)$ and $T_A=xAh$, give bounded positive units
\[
 (T_C)_0=(T_A)_0=\frac{K^2}{4D_K},\qquad
 (D/x)_0=\frac1{4K^2D_K},\quad
 (\alpha/x)_0=\frac{K^2}{4D_K}.
\]
They hold even if $h$ or $b-h$ is infinitesimal.  Specifically,
\[
 xC(1-\eta)=T_C/b,\quad xA\eta=T_A/b,\quad
 xC(1-\eta)^2=T_C(b-h)/b^2,\quad
 xA\eta^2=T_Ah/b^2.                                  \tag{JC1}
\]
Put $H_c=x(C+A)$.  This quantity is bounded below by a positive
real, although it can be unbounded.  Its reciprocal and reciprocal
square root are bounded pure coefficient elements.

\paragraph{Center the excess positive packets before taking normals.}
Retain exactly the original amplitudes
\[
 M=vW-X+uvY-uZ,\qquad J=uW+uvX-Y-vZ,
 \qquad D_u=1+u^2,\quad D_v=1+v^2.
\]
Let $F$ denote the full original parent, including its signed core
$H$.  Define $G$ by the following exact polynomial identity:
\[
 F=G+4nD_uD_v(X+ruZ)^2
             +4(m-1)(M-guZ)^2.                       \tag{JC2}
\]
An explicit expression, requiring no elimination or mixed inverse, is
\begin{align*}
G={}&4M^2+4CD_v(uW-Y)^2+4AD_uD_vY^2
       +4DJ^2+4\alpha D_u(vY-Z)^2+H\\
 &-8nrD_uD_vuZX-4nr^2D_uD_vu^2Z^2\\
 &+8(m-1)g\,uZM-4(m-1)g^2u^2Z^2.                  \tag{JC3}
\end{align*}
Every coefficient in (JC3) except $C,A$ is bounded.  At the exact
original contact
\[
 (u,v;W,X,Y,Z)=(\rho,k;\kappa\rho,-r\rho,xh/k,1)
\]
both subtracted squares and all of their first derivatives vanish.
Consequently $G$ has exactly the same zero and full stationary
contact as $F$.  This statement is essential: merely capping the
uncentered original packets would not preserve the contact.

On the finite $u=0$ residue axis use
\begin{gather*}
 X_g=-ruZ,\qquad Y_g=\eta uW,\qquad E_X=\sqrt n(X+ruZ),\\
 E_Y=\sqrt{C+A}(Y-\eta uW),\qquad
 E_M=\sqrt{m-1}(M-guZ).                              \tag{JC4}
\end{gather*}
Expand only $G$ in $E_X,E_Y$ and retain $4E_M^2$ as written.
The root $E_M$ is allowed to be redundant.  No substitution for
$W$ is made.  In particular, the positive $C$ packet introduces
no coefficient of the form $Cu^2/n$ or $Cu^2/m$.

\paragraph{The complete graph scalar and its global leading inequality.}
Write $f=G|_g=f_{ww}W^2+f_{wz}WZ+f_{zz}Z^2$.  The same exact monic
divisions as in the prior single-angle proof give
\[
 f=aP^2+b_sB^2+c_sV^2+eE^2+g_sBV,
 \qquad (P,B,V,E)=(uW/\rho,uZ,vW,\rho Z),              \tag{JC5}
\]
where
\[
 a=x\frac{f_{ww}-f_{ww}|_{u=0}}{u^2},\quad
 b_s=\frac{f_{zz}-f_{zz}|_{u=0}}{u^2},\quad
 c_s=\frac{f_{ww}|_{u=0}}{v^2},\quad
 e=\frac{f_{zz}|_{u=0}}x,\quad g_s=\frac{f_{wz}}{uv}.
\]
These are polynomial coquotients, not inverses of testing
infinitesimals.  Their coefficients are bounded by (JC1)--(JC3).
The exact constant $f_{zz}|_{u=0}=4Dv^2+4\alpha$ supplies the $x$
factor.  For finite $v$ their residues are
\[
 q(v)=\frac{K^2(1+v^2)}{D_K},\qquad
 b_s=4,\quad c_s=4,\quad g_s=-10,\qquad
 e(v)=\frac{v^2/K^2+K^2}{D_K}.                        \tag{JC6}
\]
Thus the leading scalar, using the exact identity $BV=vPE$, is
\[
                  qP^2-2vPE+eE^2+4(V-B)^2.           \tag{JC7}
\]
The key identity is now an identity for the graph of a full
bounded remainder, rather than for the overconstrained graph from
the earlier frontier note:
\[
                   q(v)e(v)-v^2
                    =\frac{(v^2-K^2)^2}{D_K^2}.      \tag{JC8}
\]
If $v$ is a finite nonzero residue with $v^2\ne K^2$, then
$10|v|<8|v|+2\sqrt{qe}$.  Splitting the exact $BV=vPE$ cross
term with the pure testing function
\[
             \theta=\frac{4|v|}{4|v|+\sqrt{q(v)e(v)}}
\]
gives strictly positive residue blocks on $(B,V)$ and $(P,E)$.
At $v_0=0$, allocate the entire cross term using $BV=vPE$;
(JC5) then has strictly positive diagonal residue
$(q(0),4,4,e(0))$ and zero cross residues.  There is no exceptional
$n_0=0$ corner in the joint regime.

Here is also a global equality and Hessian check needed at the
matched angles.  For $z=v^2>0$ and $s>0$, set
\[
 \mathcal D(s,z)=(q s+4z)(4s+e)-25zs.
\]
Exactly,
\begin{align*}
 \frac{\mathcal D(s,z)}s
 &=4qs+4ez/s+qe-9z\\
 &=4\bigl(\sqrt{qs}-\sqrt{ez/s}\bigr)^2
       +(\sqrt{qe}-\sqrt z)(\sqrt{qe}+9\sqrt z).    \tag{JC9}
\end{align*}
The second term is nonnegative by (JC8); its equality forces
$z=K^2$.  The first then forces $s=1$.  Moreover the second term is
\[
 \frac{(z-K^2)^2}{D_K^2}
       \frac{\sqrt{qe}+9\sqrt z}{\sqrt{qe}+\sqrt z},
\]
whose quadratic coefficient at contact is strictly positive.
The first square supplies a nonzero $s$ derivative there.
Hence $\mathcal D$ has a positive definite Hessian at $(1,K^2)$.
These radicals are used only in the residue-field inequality; no
radical of an element mixing coefficient and testing fields is used.

\paragraph{Restore every normal, with bounded coefficients.}
From (JC2)--(JC4) the entire parent is exactly
\[
 F=f+E_XL_X+E_YL_Y+
 H_{XX}E_X^2+2H_{XY}E_XE_Y+H_{YY}E_Y^2+4E_M^2.        \tag{JC10}
\]
The first normals are $L_X=G_X|_g/\sqrt n$ and
$L_Y=G_Y|_g/\sqrt{C+A}$.  The only unbounded weights $C,A$ do not
enter the $X$ row.  Thus
\[
 L_X=V\ell_w+B\ell_z
\]
with bounded coefficient functions carrying the infinitesimal
pure coefficient factor $1/\sqrt n$.  Both have identically zero
residue on the entire finite $u=0$ axis, including $v=0$.

For the second normal the exact form is
\[
 L_Y=\frac1{\sqrt{H_c}}\left\{
 P(B_2+xe_w)+B(\rho u e_z)+E e_*\right\},            \tag{JC11}
\]
where
\[
 B_2=8D_v\{(xA)\eta D_u-(xC)(1-\eta)\},\qquad
 e_*=8v(D-\alpha).
\]
All $e_w,e_z$ are bounded: all new terms in (JC3) have bounded
coefficients, and their $Y$ derivative has an additional $u^2Z$
factor.  Equation (JC1) makes $B_2$ bounded with identically zero
residue, even at either closing complementary opening.  Thus every
coefficient in (JC11) has zero residue.

The entire hard block is bounded and has residue
\[
             \operatorname{diag}(4D_{v,0},4D_{v,0},4)
                   \quad\hbox{on }(E_X,E_Y,E_M).      \tag{JC12}
\]
The first entry comes from the retained centered $n$ packet;
the bounded $G$ contribution is divided by $n$.  The second is
\[
 4D_v\frac{C+AD_u}{C+A}
   =4D_v\left(1+\frac A{C+A}u^2\right)
\]
plus bounded terms divided by $C+A$.  The mixed entry contains
$1/(\sqrt n\sqrt{C+A})$ times a bounded function.  The third
entry is the unexpanded centered $m$ packet.  This proves the
full-parent assertion at all remote finite-axis tests, by the
strict Gram in (JC5) and the zero-residue first normals.

\paragraph{The matched finite-angle residues.}
At $v_0=K$, normalize $Z=1$ and suppose $W$ is infinitesimal.
Use exactly the prior contact conductors
\[
 R=(uW-\kappa x)/\rho,\quad S=u-\kappa^{-1}W,
 \quad U=\rho(v-k),\quad T=u(v-k),
\]
with exact identities
\[
 u^2-x=uS+\kappa^{-1}\rho R,\quad
 W^2-\kappa^2x=-\kappa WS+\kappa\rho R,
 \quad uW-\kappa x=\rho R.                           \tag{JC13}
\]
Here $f=p^2Q_b/x+B_b$, with bounded $B_b$ and
\[
 Q_b=4D_v\{(xC)(1-\eta)^2+(xA)\eta^2D_u\}.
\]
Boundedness follows from (JC1); the remaining terms use only
the bounded coefficients of (JC3).  The full exact stationary
contact of $G$ proved after (JC3) supplies every vanishing required
by the prior monic-contact-division argument.  Successive monic
differences in $u^2-x,W^2-\kappa^2x,uW-\kappa x,v-k$ therefore
give a bounded Gram on $R,S,U,T$.

For precision, its strict residue follows from the same explicit
Schur-complement step in that argument, now with
\[
 \Phi=q(v)s^2w^2+4s^2+4v^2w^2-10vsw+e(v).
\]
Completing the $w$ square gives
\[
 \Phi=(qs^2+4v^2)
       \left(w-\frac{5vs}{qs^2+4v^2}\right)^2
       +\frac{\mathcal D(s^2,v^2)}{qs^2+4v^2}.
\]
Equation (JC9) proves a positive definite three-variable Hessian
at $(s,w,v)=(1,1/K,K)$.  Consequently the residue before the
conductor-zero correction has form
\[
 q(K)R^2+4S^2+A_1RU+B_1ST+E_1U^2,
 \qquad E_1>A_1^2/(4q(K))+B_1^2/16.
\]
Choose a real $\delta$ strictly between $B_1^2/16$ and
$E_1-A_1^2/(4q(K))$, and add $\delta$ times the exact zero
\[
 T^2-U^2-ST(v-k)-\kappa^{-1}RU(v-k)=0.
\]
The two residue blocks $(R,U)$ and $(S,T)$ are strictly positive.

An explicit equivalent residue check can be used in place of this
Schur choice.  The capped current has the exact conductor lift
\[
 (M-gu)|_g=-v\kappa S+\kappa T
                  +\eta uv\rho R+\eta u\kappa\rho U.
\]
Retaining its $T^2$ term, rather than first moving that term to $U^2$,
gives the scalar residue
\[
 K^2R^2+\frac{2(K^2-1)}{D_K}RU+\frac{U^2}{K^2}
                  +4(S-T/K)^2.
\]
Adding $2/D_K^2$ times the same conductor zero leaves the $(R,U)$
Schur margin $2/D_K^2>0$, and leaves
$4(S-T/K)^2+(2/D_K^2)T^2$ on $(S,T)$, also strictly positive.

The full-normal lifts are also the ones actually required for
(JC10), rather than scalar-only lifts.  Write
$L_X=uA_1'+WB_1'=-\kappa B_1'S+u(A_1'+\kappa B_1')$.
Both functions and all their bounded monic coquotients retain
$1/\sqrt n$.  Exact $G_X=0$ at contact supplies vanishing of
$A_1'+\kappa B_1'$.  Using (JC13) gives zero-residue coefficients
on $R,S,U,T$.  For (JC11), write its numerator as $RB_2+\rho E_2$.
Both $B_2,E_2$ have identically zero leading functions and $E_2$
vanishes at contact.  The same monic divisions give factors
$x,\rho u,\rho W$, or a zero-leading-function $(v-k)$ coquotient;
division by $\sqrt{H_c}$ preserves boundedness and zero residue.
Thus the complete seven-root Gram on
$(R,S,U,T,E_X,E_Y,E_M)$ is strict in residue.

If $W/Z$ is not infinitesimal, normalize $W=1$ and keep $Z$
bounded.  Then the scalar has a strict residue on $(P,1)$:
its positive constant contribution is $4K^2$.  All normalized
first normals have zero residue as above.  For $Z=0$, the original
parent is already the literal sum of its event packets and three
positive core squares.  The simultaneous sign copy gives $v_0=-K$.

\paragraph{Every other independent test.}
The outer-chart proof need not use the centered decomposition.
Use the original parent, cap $C,A$ at $\sigma/x$, where
$\sigma=\min(1,xC,xA)$ is a positive bounded unit, and cap $m$ at
$1$.  Retain the original nonnegative excess packets and the
entire original $n$ packet.  The four exact outer-core identities
in the prior single-angle theorem apply without change: their
strict residues require only $\sigma_0>0$, the selected $m=1$,
and $(D/x)_0>0$.  They never bound $n$.  This covers the finite
unit $u$ axis, both reciprocal axes, all their intersections, and
literal projective ends.  The inherited physical
noncommuting-residue theorem covers the remaining tests.

These charts and the entire finite $u=0$ axis exhaust all independent
projective measurement pairs.  All state charts were restored above;
complex states split into their real and imaginary parts as in the
prior proof.  Every inverse or radical used in a tensor coefficient
is pure coefficient or pure testing.  No comparison of a source
infinitesimal with a testing infinitesimal, and no mixed inverse
or radical, occurs.  Restoring the retained packets proves the theorem.
